% TOP_PROBLEM: {"id":30007537,"problem_number":"LOCAL-30007537","title":"Restructuring with diverse creditors","category_id":21,"category":{"id":21,"name":"economics","display_name":"Economics","description":"Open economic theory statements, published research agendas, and proposed scoped specifications; question provenance and historical priority are qualified per record.","slug":"economics","order_index":21,"created_at":"2026-10-08T00:00:00Z"},"status":"research_question","external_url":null,"legacy_ids":[],"published":false,"statement":"Design a feasible sovereign restructuring mechanism for heterogeneous creditors, including participation, private information, holdouts, and effects on initial lending.","economics":{"original_id":26,"field":"Fiscal policy and sovereign debt","kind":"design","formalization_basis":"adapted_research_question","source_status":"research_agenda","priority_status":"earliest_found","priority_note":"This is the earliest explicit relevant statement verified in this audit, not a claim of original priority; older literature and earlier versions may contain antecedents.","earliest_verified_reference":{"authors":["Kris James Mitchener","Christoph Trebesch"],"title":"Sovereign Debt in the 21st Century","year":2021,"url":"https://www.ifo.de/DocDL/cesifo1_wp8959.pdf","doi":"10.3386/w28598","locator":"Section 8, p. 46, final discussion of official sovereign lending; Section 10, pp. 51–52","evidence_summary":"The survey says understanding of the rise of official creditors and its implications for governments and private investors is only beginning. The restructuring mechanism below narrows this agenda."},"current_reference":{"authors":["Kris James Mitchener","Christoph Trebesch"],"title":"Sovereign Debt in the 21st Century","year":2021,"url":"https://www.ifo.de/DocDL/cesifo1_wp8959.pdf","doi":"10.3386/w28598","locator":"Section 8, p. 46, final discussion of official sovereign lending; Section 10, pp. 51–52","evidence_summary":"The survey says understanding of the rise of official creditors and its implications for governments and private investors is only beginning. The restructuring mechanism below narrows this agenda."},"formalization_note":"The survey explicitly identifies growing official-creditor complexity as insufficiently understood. The constrained restructuring mechanism is a scoped adaptation, not an attributed canonical theorem."},"statusQualification":"Published question or research agenda verified at the cited locator. The mathematical specification is adapted for this catalogue and includes additional assumptions; source publication does not establish first-ever priority or certify this exact model remains unsolved. The survey explicitly identifies growing official-creditor complexity as insufficiently understood. The constrained restructuring mechanism is a scoped adaptation, not an attributed canonical theorem.","statusReviewedAt":"2026-10-08"}
% FIRST_POSTED: 2026-10-08
% ENTRY_KIND: statement_only
% !TeX program = lualatex
\documentclass[11pt]{article}
\usepackage{fontspec}
\usepackage[a4paper,margin=25mm]{geometry}
\usepackage{amsmath,amssymb,mathtools}
\usepackage{xurl}
\usepackage[hidelinks]{hyperref}
\newcommand{\sref}[2]{\hyperlink{source:#1:#2}{[S#2]}}
\setlength{\emergencystretch}{3em}
\begin{document}
% problemId:30007537
\section*{Restructuring with diverse creditors}\label{30007537}
\subsection{Definitions and mathematical statement}
Fix an insolvent sovereign owing claims \(b_j>0\) to a finite set of creditors \(j\), including domestic banks, private foreign investors, multilateral institutions, and official bilateral lenders. Each creditor has a declared seniority, outside option, legal enforcement capacity, and participation constraint. The sovereign has stochastic repayment capacity \(R_t\) and a real social cost from delayed settlement. A restructuring mechanism chooses repayments \(x_{jt}\), creditor voting rules, and access to new financing after reports about capacity and creditor types.
Find mechanisms minimizing expected discounted crisis costs subject to aggregate feasibility
\[
\sum_j x_{jt}\le R_t+\text{permitted new financing at }t,
\]
creditor participation, any statutory seniority constraints, and incentive compatibility for specified private information. Characterize when participation by diverse official and private creditors permits timely settlement, and when holdouts or incompatible legal claims make the target allocation unattainable. The mechanism must include how anticipated restructuring terms affect initial lending and borrowing; a one-period negotiated haircut is insufficient. Identify empirical moments that discriminate between legal, informational, and coordination barriers using documented claim structures and settlement histories. This is a restricted mechanism-design target for a stated creditor environment, not a claim that all sovereign restructurings admit one universally optimal procedure.

\par\textbf{Formulation and scope.} The survey explicitly identifies growing official-creditor complexity as insufficiently understood. The constrained restructuring mechanism is a scoped adaptation, not an attributed canonical theorem.

\par\textbf{Open-question provenance.} An explicit question or research agenda was verified at the source locator. This does not by itself certify that every later version remains unresolved \sref{30007537}{1}.

\par\textbf{Historical priority.} This is the earliest explicit relevant statement verified in this audit, not a claim of original priority; older literature and earlier versions may contain antecedents.

\subsection{Short English statement}
Design a feasible sovereign restructuring mechanism for heterogeneous creditors, including participation, private information, holdouts, and effects on initial lending.

\subsection{Sources}
\begin{itemize}\raggedright
\item[\textnormal{[S1]}]\hypertarget{source:30007537:1}{} Kris James Mitchener, Christoph Trebesch. 2021. Sovereign Debt in the 21st Century. Section 8, p. 46, final discussion of official sovereign lending; Section 10, pp. 51–52.
\url{https://www.ifo.de/DocDL/cesifo1_wp8959.pdf}.
DOI: 10.3386/w28598.
{\small\textit{Earliest verified question/agenda reference; historical priority qualified below. The survey says understanding of the rise of official creditors and its implications for governments and private investors is only beginning. The restructuring mechanism below narrows this agenda.}}
\end{itemize}
\end{document}
